\section{Explicit covering calculations and derived transport}
\label{cover:details}
Here are the linear details behind Proposition~\ref{cover:band}. Write
$x_i=F_{ii}$ for diagonal entries of an endomorphism $F$ preserving the
three vertex spaces. Solving $FX=XF$, $FY=YF$ leaves a common diagonal
$c$ and six free entries. More explicitly its only off-diagonal entries are
\[
 F_{02}=F_{13}=t_0,\quad F_{15}=t_1,\quad
 F_{35}=F_{46}=t_2,\quad F_{48}=t_3,\quad
 F_{68}=F_{70}=t_4,\quad F_{72}=t_5.
\]
All others vanish. This can be checked column by column using the displayed
partial permutation matrices $X,Y$: their nonzero columns specify all
constraints, and the zero columns force the remaining entries to zero.
If $r_0,\ldots,r_5$ are the six radical generators in the displayed order,
the only nonzero products are
\[
 r_0r_2=r_1,\qquad r_2r_4=r_3,\qquad r_4r_0=r_5.
\]
Every product of three radical generators vanishes. Thus the listed space
is a local algebra, not merely a seven-dimensional vector-space count.

For clarity, the Euler equations used in the obstruction are
\begin{align*}
 d_1-d_0&=1,&d_3-d_2&=1,&d_4-d_3&=1,\\
 d_6-d_5&=1,&d_7-d_6&=1,&d_0-d_8&=1,\\
 d_1-d_2&=1,&d_4-d_5&=1,&d_7-d_8&=1,
\end{align*}
where $d_i=D_{ii}$. Subtract the last row from the first two rows after
summing: the left side is zero and the right side is three. This is an
identity in every characteristic-zero coefficient field. In characteristic
three this particular obstruction disappears; no conclusion there is
asserted. The closed walk
\[
 a_0b_0^{-1}a_0\,a_1b_1^{-1}a_1\,a_2b_2^{-1}a_2
\]
explains the computation: six forward steps and three backward steps have
net lift displacement three. Arbitrary changes of basis cannot remove the
Euler contradiction. This is an algebraic lift obstruction, not an
assertion about the topological fundamental group of the compact torus.

The Hom formula also has a constructive perfect-complex version. Each
cover representable is finite support and pushes to its corresponding
vertex projective. For a bounded complex $P^\bullet$ of finite sums of such
representables and a bounded finite-support complex $N^\bullet$, decompose
the entire Hom complex by deck displacement. The differential preserves
that displacement, and only finitely many displacements occur across the
finitely many degrees. Taking cohomology gives
\[
 \operatorname{Hom}_{D(B_L)}(FP^\bullet,FN^\bullet[q])
 \simeq\bigoplus_g
 \operatorname{Hom}_{D(\widetilde B_L)}(P^\bullet,T_gN^\bullet[q]).
\]
Bounded projective complexes compute these derived Hom spaces, which
justifies the equality rather than assuming every module perfect. This
elementary proof does not assert an equivalence of entire bounded derived
categories or cover objects of infinite projective dimension. The broader
precovering statement has its exact hypotheses in
\cite[Theorem~4.28]{CoveringLecture2026}.

Finally the branch involution $\sigma(a_i)=b_i$, $\sigma(b_i)=a_i$
preserves $B_L$. In characteristic different from two one can study
$B_L\#k[\Z/2]$ through equivariant representations. For the nine-dimensional
module above, $\operatorname{rank}X=6$ and $\operatorname{rank}Y=3$, so
$M\not\simeq\sigma M$. The sum $M\oplus\sigma M$ has canonical swap
equivariance. This is an explicit orbit construction; it neither chooses a
nonexistent fixed-object linearization nor identifies this skew group
algebra with the particular GLS algebras in \cite{TameGLS2026}.

\texttt{CoveringBand.lean} verifies the actual matrices' mixed products,
cubic nilpotence and absence of an arbitrary Euler matrix. The general
covering, module endomorphism classification, derived functors and the
surface/Fukaya equivalence are not fully formalized by these three matrix
statements. Exact rational code additionally checks thirty-six Hom pairs,
the full radical multiplication and the two Hom-control defects.
